This paper develops an aesthetics of the everyday and, on its basis, a theory of aesthetic education. Its guiding claim is that everyday aesthetic perception, small in amplitude and prone to go dead through habituation, is generated rather than given, coming to be in the relation between a perceiver and what is perceived and, in the cases that matter most, between two perceivers turned together toward a shared world. From this it draws its central proposal, \emph{generative relational aesthetic education}, as one thesis with four faces: that beauty is generated within a relational structure; that an education in beauty must therefore cultivate the relation and not only the individual, which marks it off from the individualist tradition from Schiller onward; that the relational and the individual stand in a spiral dialectic in which neither is prior; and that aesthetic generation carries an ethical valence at its root, so that beauty and ethics share a single source rather than being joined by a later moral rule.

The argument proceeds in three spines, of generation, of the body, and of justice, arguing across phenomenology, cognitive science, psychoanalysis, and the political economy of the senses that perception constructs a world under a prior frame; giving this its empirical ground in the neuroscience of perceptual plasticity; and showing, through the distribution of the sensible, that the uneven distribution of perceptual capacity is a neglected injustice. On this ground the paper offers, as a preliminary and openly revisable contribution, four models of everyday aesthetic practice, each with its theoretical descent made explicit: a basic generative cycle; a reflective variant beginning in silence, where reception suspends private judgement; the leaving of emptiness, a field in which contingency may generate of itself; and an iterative symbolic refinement modelled on evolution without a fitness function. The paper closes not in synthesis but in directed fissures, the first pointing toward the emergence of culture in intimate relation.
